\usepackage[a4paper,margin=2cm]{geometry}
\usepackage{xcolor}
\usepackage{tcolorbox}
\tcbuselibrary{raster}
\usepackage{enumitem}

\definecolor{murfyBlue}{HTML}{0081FB}
\definecolor{murfyBlueDark}{HTML}{004D9F}
\definecolor{murfyBlueSoft}{HTML}{E6F5FF}
\definecolor{murfyInk}{HTML}{212528}
\definecolor{murfyMuted}{HTML}{868D94}
\definecolor{murfyLine}{HTML}{CDD3D8}

\newcommand{\murfyTitleBlock}[3]{%
  {\color{murfyInk}\fontsize{25}{29}\selectfont\bfseries #1}\par\vspace{5pt}
  {\color{murfyMuted}\large #2}\par\vspace{3pt}
  {\color{murfyMuted}\small #3}\par\vspace{10pt}
  {\color{murfyBlue}\rule{\linewidth}{1.6pt}}\par
}
\newtcolorbox{murfyCard}{
  colback=white, colframe=murfyLine, arc=3pt, boxrule=0.4pt,
  left=11pt, right=11pt, top=9pt, bottom=9pt,
  before upper={\raggedright\hyphenpenalty=10000\relax},
}
\newtcolorbox{murfyNote}{
  colback=murfyBlueSoft, colframe=murfyBlueSoft, arc=3pt,
  left=13pt, right=13pt, top=10pt, bottom=10pt, boxrule=0pt,
}

\newcommand{\murfyCardTitle}[1]{{\color{murfyBlueDark}\bfseries #1}\par\vspace{3pt}}
\newcommand{\murfySectionTitle}[1]{{\color{murfyInk}\large\bfseries #1}\par\vspace{7pt}}
\newcommand{\murfyBulletLabel}{\textcolor{murfyBlue}{\textbullet}}

% Murfy Guide
%
% This file only draws the guide page. It is not needed to write your paper:
% delete the two guide lines in main.tex and this file to remove it.

\newcommand{\murfyguide}{%
  \thispagestyle{empty}%
  \begingroup
  \setlength{\parindent}{0pt}
  \sffamily
  \color{murfyInk}

  \murfyTitleBlock
    {Murfy Guide}
    {The LaTeX editor with AI built in}
    {Delete this page once you start your paper.}

  \vspace{16pt}
  \murfySectionTitle{Your first PDF in three steps}
  \begin{enumerate}[leftmargin=22pt, itemsep=5pt, topsep=0pt,
                    label=\textcolor{murfyBlue}{\bfseries\arabic*.}]
    \item \textbf{Write.} Add your text under \texttt{\textbackslash section} in
          \texttt{main.tex}.
    \item \textbf{Compile.} Your edits are typeset into the PDF. If something breaks, the fix
          button in the compile log lets AI explain and repair it.
    \item \textbf{Remove this guide.} Delete the \texttt{\textbackslash murfyguide} line in
          \texttt{main.tex} and this page is gone.
  \end{enumerate}

  \vspace{13pt}
  \murfySectionTitle{What you can do here}
  \begin{tcbraster}[raster columns=2, raster equal height=rows, raster valign=top,
                    raster column skip=9pt, raster row skip=9pt]
    \begin{murfyCard}
      \murfyCardTitle{AI paper review}
      Get an AI review before you submit. Murfy reads your paper like a reviewer and reports the
      summary, strengths and weaknesses.
    \end{murfyCard}
    \begin{murfyCard}
      \murfyCardTitle{AI compile error fix}
      No need to read the log yourself. Murfy points to the failing line, explains the cause and
      applies the fix.
    \end{murfyCard}
    \begin{murfyCard}
      \murfyCardTitle{Files and formats}
      Add \texttt{.tex}, \texttt{.bib} and image files as you go. Need a journal format? Create a
      project from a template on the dashboard.
    \end{murfyCard}
    \begin{murfyCard}
      \murfyCardTitle{Write together}
      Invite co-authors and edit the same document in real time. Leave comments on a sentence to
      discuss it in place.
    \end{murfyCard}
  \end{tcbraster}

  \vspace{13pt}
  \murfySectionTitle{Good to know}
  \begin{itemize}[leftmargin=18pt, itemsep=5pt, topsep=0pt, label=\murfyBulletLabel]
    \item Project settings let you pick the TeX Live version and compiler your journal asks for.
    \item Your edit history is kept, so you can roll back or download an earlier state.
    \item Invite links come in two kinds: one for editing, one for viewing only.
  \end{itemize}

  \vspace{13pt}
  \begin{murfyNote}
    {\bfseries To remove this guide page}\par\vspace{3pt}
    Delete the \texttt{\textbackslash input\{murfy-guide.tex\}} and
    \texttt{\textbackslash murfyguide} lines from \texttt{main.tex}, then delete the
    \texttt{murfy-guide.tex} file. Your paper is not affected.
  \end{murfyNote}

  \endgroup
  \clearpage
}
